\cleardoublepage
\thispagestyle{plain}

\begin{center}
    {\Large \bfseries FICHA RESUMEN}
\end{center}

\vspace{0.3cm}

\noindent El presente proyecto de grado aborda la generación automática de la capa de persistencia en aplicaciones mediante el desarrollo de un módulo traductor de código enfocado en Django ORM. La propuesta integra la Gramática de Idioms de un diagramador entidad-relación y técnicas de procesamiento de lenguaje natural aplicadas a código mediante el modelo de lenguaje CodeT5.

\vspace{0.3cm}

\noindent A través de una metodología iterativa basada en prototipos, se diseñó e implementó un pipeline que procesa la representación semántica derivada del análisis sintáctico (mediante ANTLR) y la transforma en declaraciones de modelos Django robustas y funcionales. El sistema fue evaluado bajo escenarios de generación zero-shot y fine-tuning, alcanzando métricas óptimas de exactitud sintáctica y semántica en la traducción.

\vspace{0.8cm}

% Palabras Clave
\noindent \textbf{Palabras Clave:} Django ORM, CodeT5, ANTLR, Gramática de Idioms, Generación de Código, Capa de Persistencia, Modelos de Lenguaje.